\section{Background: Anatomy of a Harness}\label{sec:background}
\answers{the descriptive part of RQ1 and sets up RQ2}

\paragraph*{The loop} A harness builds a context, calls the model, checks each proposed tool call against a permission policy, executes it and feeds the result back; repository content, tool output and MCP results flow in, commands and requests flow out.

\paragraph*{Repository files} Harnesses auto-load instruction files (\texttt{AGENTS.md}, \texttt{CLAUDE.md}, \texttt{.github/copilot-instructions.md}, \texttt{.cursor/rules}), settings that can set environment variables, permission modes and \emph{hooks} (shell commands bound to lifecycle events), and MCP declarations (\texttt{.mcp.json}, \texttt{.vscode/mcp.json}). Formats are shared: Grok Build imports Claude Code and Cursor hook and MCP files\evD~\cite{grokbuild_hooks,grokbuild_mcp}. A cloned repository is therefore both data and configuration.

\paragraph*{Permissions and sandbox} Modes range from ``ask before every tool call'' to ``approve everything'' (YOLO, bypass, full-auto, always-approve), with allow and deny rules over commands and paths. Several harnesses now default to a model-based approver; Claude Code starts \texttt{claude -p} in a classifier-mediated \texttt{auto} mode when feature flags are not fetched\evD\evT~\cite{cc_permission_modes}. Some confine commands with an OS sandbox (Seatbelt, bubblewrap, Landlock) or a container\evD~\cite{cc_sandboxing}. Approval decides \emph{whether} an action runs; the sandbox limits \emph{what it reaches}.

\paragraph*{Extensions, state and headless mode} MCP~\cite{mcpspec} servers are launched locally (often via unpinned \texttt{npx} or \texttt{uvx}) or reached remotely, and their tool descriptions enter the context verbatim; skills and plugins add code and prompts from marketplaces. Credentials, transcripts and persistent memory live in a per-user state directory, and sub-agents may inherit the parent's permissions. The harnesses we benchmarked all run non-interactively\evD\evT{} (\texttt{claude -p}, \texttt{codex exec}, \texttt{grok -p}, \texttt{copilot -p}); in CI the triggering issue or pull request is attacker-controllable and the job may hold secrets.

\paragraph*{The 2026 landscape} The field splits into open-source harnesses, closed but installable clients, hosted agents and SDK runtimes (Section~\ref{sec:method}). On 31 March 2026 the npm release 2.1.88 of Claude Code was reported to ship a source map exposing client internals\evTP~\cite{cc_sourcemap2026}; we cite this only as an event and did not use the exposed material.
